\section{从 EVT 到量产：验证、EMC、良率与失效分析}

一块工程样机完成受限流首轮上电，只证明“这一块板在当前条件下走通了启动路径”。产品要进入量产，还要证明设计在电压、温度、负载、接口兼容和电磁环境的规定范围内满足需求，生产线能稳定复制它，并且失败样品能沿证据链回到设计、物料、工艺或测试程序。EVT、DVT、PVT 是业界常用的阶段名称，但不同企业对样本数、门槛和职责的定义并不完全相同；真正有约束力的是本项目受控的需求、验证计划和放行标准。

\topic{阶段门不是日历节点，而是一组可审核证据}

EVT（engineering validation test）主要验证工程原理和电气实现：电源、时钟、复位、DDR/高速链路、热路径和关键功能是否在设计角落仍有裕量。DVT（design validation test）把对象扩展到接近产品形态的硬件、固件和结构件，覆盖法规/行业符合性、环境、可靠性、兼容性和用户场景。PVT（production validation test）用接近量产的物料、工装、测试站和作业流程验证节拍、测试覆盖与良率。随后才是受控量产爬坡与正式量产。

\begin{pdfdiagram}[x=1cm,y=1cm]
  \node[hw state, minimum width=2.35cm] (req) at (1.25,4.9) {需求与风险\\可测验收条件};
  \node[hw control, minimum width=2.35cm] (evt) at (4.1,4.9) {EVT\\原理与电气裕量};
  \node[hw control, minimum width=2.35cm] (dvt) at (6.95,4.9) {DVT\\产品与符合性};
  \node[hw control, minimum width=2.35cm] (pvt) at (9.8,4.9) {PVT\\产线与测试复制};

  \node[hw state, minimum width=2.35cm] (ramp) at (9.8,2.45) {量产爬坡\\版本、节拍、良率};
  \node[hw compute, minimum width=2.35cm] (telemetry) at (6.95,2.45) {生产遥测\\FPY、缺陷与返修};
  \node[hw control, minimum width=2.35cm] (fa) at (4.1,2.45) {失效分析\\复现与根因证据};
  \node[hw state, minimum width=2.35cm] (change) at (1.25,2.45) {纠正措施\\设计/工艺/测试};

  \node[hw control, minimum width=2.35cm] (verify) at (1.25,0.0) {验证修正\\新批次与回归};
  \node[hw state, minimum width=2.35cm] (release) at (4.1,0.0) {新 release\\散列值与生效点};
  \node[hw compute, minimum width=2.35cm] (monitor) at (6.95,0.0) {持续监控\\趋势与预警线};
  \node[hw state, minimum width=2.35cm] (closed) at (9.8,0.0) {问题关闭\\证据可追溯};

  \draw[hw bus] (req) -- (evt);
  \draw[hw bus] (evt) -- (dvt);
  \draw[hw bus] (dvt) -- (pvt);
  \draw[hw bus] (pvt) -- (ramp);
  \draw[hw bus] (ramp) -- (telemetry);
  \draw[hw bus] (telemetry) -- (fa);
  \draw[hw bus] (fa) -- (change);
  \draw[hw bus] (change) -- (verify);
  \draw[hw bus] (verify) -- (release);
  \draw[hw bus] (release) -- (monitor);
  \draw[hw bus] (monitor) -- (closed);
\end{pdfdiagram}
\diagramnote{从需求到量产问题关闭的三层闭环。EVT、DVT、PVT 是逐步扩大的验证阶段；量产数据又把真实缺陷送回失效分析。临时返修只有进入受控设计、工艺或测试 release，并由新批次验证后，才成为可复制的纠正措施。}

\begin{longtable}{L{0.14\linewidth}L{0.24\linewidth}L{0.29\linewidth}L{0.23\linewidth}}
\toprule
阶段 & 主要问题 & 典型证据 & 放行不代表 \\
\midrule
EVT & 原理图、PCB、电源/信号/热设计能否工作并有裕量 & 波形、眼图、训练窗口、温升、功能 trace、设计修改单 & 结构、法规、全场景与产线均已完成 \\\tablerowrule
DVT & 接近最终产品的设计能否满足完整需求与环境 & 需求覆盖矩阵、EMC/安全/环境报告、兼容性和可靠性结果 & 量产工装、测试节拍和供应链已稳定 \\\tablerowrule
PVT & 量产物料、工艺、工装和测试站能否稳定复制 & 工艺能力、站点 GR\&R、FPY、节拍、追溯码和小批构建报告 & 长期批次漂移不会发生 \\\tablerowrule
Ramp & 产量提升时良率、供应与维修系统能否保持受控 & 分批良率趋势、缺陷 Pareto、返修闭环、变更生效点 & 可以停止趋势监控 \\\tablerowrule
MP & 产品与过程在批准基线下持续生产 & release、偏差许可、过程控制与现场反馈 & 任意换料或工艺调整都自动等价 \\
\bottomrule
\end{longtable}

阶段之间允许有受控的未关闭问题，但必须写明风险、临时措施、责任人、截止点和不可跨越条件。把失败项改成“仅供参考”并不能降低风险；若需求本身需要修改，也要经过同样的评审和版本控制，避免通过移动验收线伪造通过。

\topic{验证矩阵把一句需求变成可复现测量}

“系统应稳定”“EMI 要低”“温度不能太高”都不能直接验收。每项需求应连接到测试对象、硬件/固件版本、仪器与校准状态、连接方式、输入激励、温压和负载角落、样本策略、通过阈值、原始数据与报告位置。这样一次失败才能复现，一次通过也能说明覆盖了什么。

\begin{longtable}{L{0.18\linewidth}L{0.28\linewidth}L{0.26\linewidth}L{0.18\linewidth}}
\toprule
矩阵字段 & 必须回答的问题 & 缺失时的风险 & 合格证据 \\
\midrule
需求与风险 ID & 为什么要测，失败影响是什么 & 测试很多却漏掉关键风险 & 可追到批准需求和风险分析 \\\tablerowrule
样机与版本 & 哪块板、何种 BOM/固件/结构件 & 不同版本结果被混在一起 & 序列号、release 与散列值 \\\tablerowrule
方法与配置 & 仪器、探头、线缆、软件和布置怎样复现 & 测量系统改变结果 & 校准信息、照片/图纸与脚本版本 \\\tablerowrule
角落与样本 & 哪些电压、温度、负载和器件分布被覆盖 & 只验证典型样机 & 预先批准的样本和角落计划 \\\tablerowrule
通过阈值 & 数值、统计或行为边界是什么 & 结果出来后主观解释 & 测试前冻结的判定规则 \\\tablerowrule
原始证据 & 波形、日志、失败品和计算过程在哪里 & 报告结论无法复核 & 受控数据、时间戳和操作者 \\
\bottomrule
\end{longtable}

测试样本数要由风险、失效分布、成本和法规要求决定，不能把一个固定数量套给所有项目。发现一次失败后，也不能只增加样本直到平均值变好；应先判断这是测量误差、个体离群、批次问题还是设计系统性问题，再决定补测、隔离或修改。

\topic{EMC 同时约束“向外发出多少”和“外界来了还能否工作”}

电磁兼容包括 emission 与 immunity 两个方向。传导/辐射发射测试关注产品通过电源线、信号线、机壳和空间向外耦合的骚扰；抗扰度测试关注静电、射频场、快速瞬变、浪涌或其他规定扰动到来时，产品是否保持规定性能或按允许方式恢复。适用项目、频段、限值、端口和性能判据由销售地区、产品类别与对应标准决定，不能从另一类产品的报告复制。\cite{iec-emc}

正式实验室测试之前应做 pre-compliance。用与最终线缆、机壳、电源和工作模式接近的样机，在可控布置中扫描高风险工作状态；近场探头用于寻找板上热点，电流探头用于观察线缆共模电流，示波器和频谱仪用于把峰值与时钟、开关电源、数据突发或刷新周期关联。预扫结果不能替代认可实验室报告，但能在修改成本较低时暴露主要耦合路径。

\begin{longtable}{L{0.18\linewidth}L{0.25\linewidth}L{0.28\linewidth}L{0.19\linewidth}}
\toprule
测试方向 & 主要观察量 & 诊断线索 & 设计控制点 \\
\midrule
传导发射 & 电源/线缆端口上的骚扰电压或电流 & 开关频率及谐波、共模回路、滤波器谐振 & 输入滤波、回路面积、接地与线缆终端 \\\tablerowrule
辐射发射 & 空间电磁场强或规定接收量 & 时钟谐波、连接器、缝隙和长线缆热点 & 回流连续性、屏蔽、边沿速度与布局 \\\tablerowrule
射频抗扰度 & 外加连续场/注入下的功能性能 & 模拟前端、复位、时钟或通信误码 & 滤波、屏蔽、阈值裕量与软件恢复 \\\tablerowrule
ESD & 接触/空气放电后的性能与损伤 & 外壳开口、按键、连接器到敏感地的路径 & 低阻放电路径、钳位位置与隔离间距 \\\tablerowrule
EFT/浪涌 & 电源或长线端口瞬态下的保持与恢复 & 保护器件能量、地弹、复位与掉电路径 & 分级保护、能量预算、滤波和 brownout 策略 \\
\bottomrule
\end{longtable}

\topic{EMC 失败要从频率峰值走回电流闭环}

看到某个超限频点后，先核对测试布置、校准和工作模式，再将频率与系统时钟、DDR 数据率、DC--DC 开关频率及其谐波相关联。随后用近场扫描定位能量源，用电流探头判断线缆是否成为天线，沿参考平面、连接器和机壳寻找共模转换路径。一次只改变一个可解释因素，例如降低可编程边沿速度、增加临时共模抑制、改善回流搭接或屏蔽一个开口，再用同一布置复测。

\begin{longtable}{L{0.10\linewidth}L{0.25\linewidth}L{0.29\linewidth}L{0.26\linewidth}}
\toprule
步骤 & 动作 & 要保存的证据 & 退出条件 \\
\midrule
E0 & 复核实验室/预扫布置与运行模式 & 照片、线缆、负载、软件和峰值列表 & 排除布置或模式误差 \\\tablerowrule
E1 & 把峰值与时钟和周期事件相关 & 基频、谐波间隔、时域触发关联 & 得到候选源而非仅有频点 \\\tablerowrule
E2 & 近场扫描板、连接器与缝隙 & 空间热点图和探头方向 & 缩小源与耦合路径 \\\tablerowrule
E3 & 测量线缆共模电流和回流连续性 & 电流谱、参考平面和机壳连接状态 & 判断主要辐射结构 \\\tablerowrule
E4 & 单变量临时修改并 A/B 复测 & 修改前后同配置曲线 & 峰值变化支持因果假设 \\\tablerowrule
E5 & 把修正进入设计 release & 原理图/PCB/BOM/结构修改与风险回归 & 正式样机重测通过 \\
\bottomrule
\end{longtable}

只在实验样机线缆上夹一个磁环可以验证共模假设，但它只有进入明确料号、装配位置和量产检查后才是产品修复。降低时钟或关闭接口也能让峰值消失，但若破坏性能需求，就只能作为定位手段。EMC 修正还要回归信号完整性、ESD 电流路径、热、成本和接口兼容，避免解决一个频点却引入新的系统故障。

\topic{最终良率可能很好，首过良率仍会揭示过程损失}

\term{首过良率}{first-pass yield, FPY} 是无需返修便通过某站或整线的良品比例：
\[
  \mathrm{FPY}=\frac{\text{首次通过数量}}{\text{进入该流程的总数量}}.
\]
多个近似独立站点的 rolled throughput yield（RTY）可用各站首过率乘积估计。若四站 FPY 依次为 98.5\%、99.0\%、99.5\% 和 98.0\%，则
\[
  \mathrm{RTY}=0.985\times0.990\times0.995\times0.980
  \approx 0.9509,
\]
也就是约 95.1\% 的投入品可在四站都不返修地通过。最终良率可能因返修升到更高，但返修工时、重复热循环、隐性损伤和交付波动仍是真实成本。

\begin{longtable}{L{0.18\linewidth}L{0.24\linewidth}L{0.27\linewidth}L{0.21\linewidth}}
\toprule
指标 & 主要回答 & 容易掩盖什么 & 合适配套视图 \\
\midrule
单站 FPY & 哪个工位首次通过能力不足 & 缺陷是否在上游产生、本站才发现 & 按站点、机台、班次和缺陷码分层 \\\tablerowrule
整线 RTY & 无返修走完整流程的概率 & 站点相关性与重复测量影响 & 流程路径和实际序列追踪 \\\tablerowrule
最终良率 & 交付前最终可接受比例 & 返修次数、成本和可靠性风险 & 返修率、报废率与维修工时 \\\tablerowrule
DPPM/缺陷率 & 单位产出中的缺陷规模 & 一个产品可能有多个缺陷 & 缺陷类型、严重度和批次趋势 \\\tablerowrule
测试逃逸率 & 坏品被测试判为通过的风险 & 尚未在现场暴露的潜在失效 & 审核/破坏性抽检与现场反馈 \\\tablerowrule
误拒率 & 良品被测试判为失败的比例 & 重测后通过造成的节拍损失 & 测量系统分析与限值裕量 \\
\bottomrule
\end{longtable}

良率必须带分层维度：板号/BOM 版本、PCB 与元件 lot、生产线、机台、治具、测试程序、操作者、班次和环境。总良率不变时，某个新料号对应的缺陷可能已快速上升，却被其他批次改善抵消。每个测试结果还要关联序列号和配置 generation，才能从 Pareto 柱状图追到可复现样品。

\topic{测试限值要在覆盖、误拒和逃逸之间建立证据}

生产测试通常比实验室验证更短，必须选择能快速暴露关键缺陷的观测点。guard band 可以在规格边界内提前拦截漂移，但过宽会提高误拒；过窄则增加逃逸。限值应来自设计裕量、测量不确定度、器件/过程分布和失效后果，不能用“把良率调到目标值”反推。

测试站本身也需要验证。治具接触电阻、探针磨损、仪器漂移、线缆和软件版本都可能制造假失败。Gauge repeatability and reproducibility（GR\&R）或等价测量系统分析，用于区分产品变差与测量系统变差。黄金样品可以监视站点漂移，但黄金样品也会老化，应有校准、保管和替换规则。

\begin{longtable}{L{0.21\linewidth}L{0.25\linewidth}L{0.28\linewidth}L{0.16\linewidth}}
\toprule
异常模式 & 区分证据 & 优先动作 & 错误反应 \\
\midrule
单站失败突然上升 & 同批产品在平行站通过，黄金样品也漂移 & 检查治具、校准和软件版本 & 立即修改产品设计 \\\tablerowrule
同一 lot 跨站失败 & 料号/PCB/工艺 lot 高度相关 & 隔离 lot，做来料与失效分析 & 反复重测直到通过 \\\tablerowrule
边界值随温度漂移 & 原始测量与环境记录相关 & 核对设计裕量和补偿策略 & 只放宽限值 \\\tablerowrule
返修后短期通过 & 返修位置与失效码相关，可靠性未知 & 验证返修工艺与应力影响 & 仅看最终良率 \\\tablerowrule
现场出现测试未见故障 & 现场模式/环境不在覆盖矩阵 & 补充可激发机制的测试与监控 & 把责任归为“用户偶发” \\
\bottomrule
\end{longtable}

\topic{失效分析先保护证据，再逐步增加侵入性}

失败品到达实验室时，先冻结序列号、板/BOM/固件版本、生产 lot、测试站日志、返修历史、使用环境和故障时间线。随后在不改变样品的条件下复现：确认电源、连接、软件和负载，记录第一条偏离正常状态的证据。互换板卡或部件可以缩小范围，但每次交换都要记录，避免把间歇故障转移后失去原始配置。

检查应从非破坏方法逐步深入：外观和显微镜、日志与边界扫描、电气静态测量、热像、示波器/协议 trace、X-ray、超声或 TDR；只有在假设和样品计划明确后，再做拆焊、染色、切片、开封等破坏分析。破坏前应保留对照品和原始影像，因为切片之后无法回到“未扰动”的焊点状态。

\begin{longtable}{L{0.10\linewidth}L{0.25\linewidth}L{0.30\linewidth}L{0.25\linewidth}}
\toprule
状态 & 主要动作 & 交付证据 & 进入下一状态的条件 \\
\midrule
F0 & 隔离失败品与同批风险品 & 序列号、lot、版本、现场与测试日志 & 证据链完整，防止继续出货 \\\tablerowrule
F1 & 在受控条件下复现并建立良品对照 & 触发条件、概率、第一偏离点 & 故障可重复或已有足够现场证据 \\\tablerowrule
F2 & 非破坏定位到电源、网络、器件或工艺区域 & 波形、热图、X-ray、协议/边界扫描结果 & 候选机制能够解释全部主要现象 \\\tablerowrule
F3 & 设计单变量验证或必要的破坏分析 & 修改前后 A/B 结果、截面/材料证据 & 根因而非仅相关性得到支持 \\\tablerowrule
F4 & 定义 containment 与永久纠正措施 & 隔离范围、临时筛选、设计/工艺/测试变更 & 风险已控制且新 release 可执行 \\\tablerowrule
F5 & 用新批次和回归验证有效性 & 良率、裕量、可靠性与副作用结果 & 修正跨目标条件稳定通过 \\\tablerowrule
F6 & 监视后续批次与现场趋势 & 生效点、趋势窗口和复发判据 & 观察期满足关闭标准 \\
\bottomrule
\end{longtable}

“更换器件后恢复”只能说明故障随器件或焊点区域移动，未必证明器件硅片本身是根因；拆焊过程可能同时重熔了裂纹焊点。“补焊后通过”也不能区分焊膏、回流、焊盘设计、板翘或机械应力。根因必须解释故障机制、批次/条件分布和所有关键证据，并能预测纠正后哪些测量会改变。

\topic{问题关闭需要把知识写回可生产对象}

containment 用于立即限制风险，例如隔离 lot、增加临时筛选、暂停某个替代料或降低出货范围；corrective action 则消除已证实根因，例如修改回流窗口、焊盘、去耦、固件时序或测试覆盖。两者必须区分，避免临时全检长期掩盖过程失控。

永久措施要进入对应受控对象：原理图/PCB、BOM/AVL、固件、结构件、工艺参数、检验规范或测试程序，并写明从哪个序列号、lot 和日期开始生效。验证既要重现原故障条件，也要回归可能受影响的性能、EMC、热、可靠性与制造性。最后用后续批次良率和现场窗口确认不再复发，问题才可以关闭。

从 EVT 到量产的完整通路因此与处理器事务一样有明确提交点：需求定义成功条件，验证阶段建立设计与符合性证据，PVT 和爬坡证明生产过程可以复制，良率遥测发现偏移，失效分析保存并收敛证据，纠正措施进入新 release，后续批次验证后才提交关闭。产品化不是硬件原理之外的附录，而是让同一套原理在成千上万块实物上继续成立的最后一层工程。
